% Resume: Arumalla Mohit Krishna
% Compile with pdfLaTeX.

\documentclass[letterpaper,10pt]{article}

\usepackage[margin=0.5in]{geometry}
\usepackage[T1]{fontenc}
\usepackage{lmodern}
\usepackage{microtype}
\usepackage{titlesec}
\usepackage{enumitem}
\usepackage[hidelinks]{hyperref}
\usepackage{tabularx}
\input{glyphtounicode}

\pdfgentounicode=1
\pagestyle{empty}
\urlstyle{same}
\raggedright
\setlength{\parindent}{0pt}
\setlength{\tabcolsep}{0pt}

\titleformat{\section}
  {\large\scshape}{}{0em}{}[\titlerule]

\titlespacing*{\section}{0pt}{6pt}{3pt}

\newcommand{\entry}[2]{%
  \noindent
  \begin{tabularx}{\linewidth}
    {@{}>{\raggedright\arraybackslash}X@{\hspace{6pt}}r@{}}
    \textbf{#1} & #2
  \end{tabularx}\par
}

\newcommand{\bullets}{
  \begin{itemize}[
    leftmargin=1.2em,
    topsep=2pt,
    itemsep=1pt,
    parsep=0pt,
    partopsep=0pt
  ]
}

\newcommand{\finishbullets}{
  \end{itemize}
}

\begin{document}

%----------HEADING----------
\begin{center}
    {\LARGE\bfseries Arumalla Mohit Krishna}\\[3pt]
    +1 (631) 352-8643 $|$
    \href{mailto:mohitkrishna.arumalla@stonybrook.edu}
    {mohitkrishna.arumalla@stonybrook.edu} $|$
    \href{https://linkedin.com/in/mohit-krishna-arumalla}
    {LinkedIn}
\end{center}
\vspace{-8pt}

%-----------EDUCATION-----------
\section{Education}

\entry
  {\textbf{Stony Brook University}}
  {Aug 2025 -- Dec 2026}

M.S. in Computer Science \hfill New York, USA\\
\textit{
Coursework: Natural Language Processing,
Introduction to Computer Vision, Vision Language Modeling
}
\vspace{3pt}

\entry
  {\textbf{Birla Institute of Technology and Science, Pilani}}
  {Aug 2019 -- 2023}

B.E. in Electronics \& Communication; GPA: 8.02/10 \hfill India
%-----------PROFESSIONAL SUMMARY-----------
\section{Professional Summary}

\textbf{Software engineer} with 2+ years of experience building distributed services, semantic search, and real-time data pipelines, currently pursuing an M.S. in Computer Science at Stony Brook University. Experienced in \textbf{data engineering, AI-enabled capabilities, product management}, and cloud infrastructure. Skilled in Python, Go, PyTorch, with strong \textbf{problem-solving, communication, and ownership} mindset aligned with building impactful business solutions.

%-----------TECHNICAL SKILLS-----------
\section{Technical Skills}

\textbf{Languages:}
Python, Go, C++, SQL, Java\\

\textbf{Data Engineering \& AI:}
PyTorch, TensorFlow/Keras, scikit-learn, NumPy, pandas,
Transformers, AI/ML pipelines, Vector Search, Data Cleansing,
Data Integration, Model Serving\\

\textbf{Data \& Systems:}
gRPC, Kafka, Apache Flink, Flink SQL, Debezium/CDC,
PostgreSQL, Redis, Elasticsearch, Airflow, BigQuery\\

\textbf{Infrastructure:}
AWS, Docker, Kubernetes, Linux, Git, CI/CD,
Istio/Envoy, New Relic, Cloud Platforms

%-----------EXPERIENCE-----------
\section{Experience}

\entry
  {\textbf{Software Engineer II} $|$ \textbf{Infra.Market}, Bangalore}
  {Apr 2024 -- Aug 2025}

\bullets

  \item Designed and built \textbf{data pipelines} with CDC using Flink, Kafka, Debezium, and PostgreSQL, achieving near real-time data synchronization with \textbf{exactly-once processing semantics}.

  \item Developed semantic product search integrated with \textbf{machine learning embeddings} and Elasticsearch for enhanced natural-language product discovery.

  \item Implemented Istio/Envoy service mesh and New Relic monitoring across microservices to support scalable and reliable cloud infrastructure.

\finishbullets

\entry
  {\textbf{Software Engineer I} $|$ \textbf{Infra.Market}, Chennai}
  {Jul 2023 -- Mar 2024}

\bullets

  \item Owned Search and Location microservices; developed \textbf{gRPC services} processing 500K+ daily requests supporting product discovery and geolocation features.

  \item Improved reliability by fixing PostgreSQL connection orchestration issues with Gorm ORM, resulting in 85\% reduction in connection pool exhaustion.

  \item Built and optimized Airflow pipelines migrating data from BigQuery to Elasticsearch, enabling efficient and automated indexing of search data.

\finishbullets

%-----------PROJECTS-----------
\section{Projects}

\entry
  {GPU-Optimized Personalized Product Search}
  {Aug 2026 -- Oct 2026}

\textit{Python, PyTorch, BGE, NVIDIA Triton, Go, gRPC}

\bullets

  \item Developed a two-tower retrieval model using \texttt{BAAI/bge-small-en-v1.5} embeddings and trained user interaction tower for \textbf{personalized semantic search}.

  \item Created GPU-accelerated inference pipeline with Triton serving and Go/gRPC gateway optimizing \textbf{low-latency retrieval} and scalability for bursty traffic.

  \item Benchmarked latency, throughput, and accuracy metrics (Recall@K, NDCG@K) and designed overload recovery methods for production readiness.

\finishbullets

\entry
  {Attention-Based Hallucination Detection in NMT}
  {Apr 2026 -- May 2026}

\textit{Python, PyTorch, mBART, Transformers, scikit-learn}

\bullets

  \item Built hallucination detection model leveraging \textbf{attention features} from mBART for English-German translation quality estimation.

  \item Engineered cross- and self-attention feature extraction and trained classifiers to improve detection precision, achieving 0.712 F1.

  \item Proposed hallucination-aware reranking integrating localization and risk scoring to improve translation reliability.

\finishbullets

\entry
  {Prompt/LoRA for Self-Supervised Learning}
  {Ongoing}

\textit{Python, PyTorch, Vision Transformers}

\bullets

  \item Researching domain-specific visual prompting and \textbf{parameter-efficient fine-tuning} methods (LoRA) in self-supervised vision transformers to improve domain transferability.

  \item Evaluated classification performance under domain shift with and without prompt-based adapters for efficient model adaptation.

\finishbullets

\end{document}